\documentclass[a4paper,12pt]{jsarticle}
\usepackage[margin=25mm]{geometry}
\usepackage{amsmath,amssymb}
\usepackage{enumitem}

\begin{document}

\begin{center}
{\LARGE 数I・A 共通テスト本番レベル（図形）解答}
\end{center}

\vspace{1em}

\section*{第1問（三角比・図形と計量）}
\begin{enumerate}[label=\arabic*.]
  \item
  \[
  \angle ATB=180^\circ-42^\circ-57^\circ=81^\circ.
  \]
  \item 正弦定理より
  \[
  \frac{AT}{\sin57^\circ}=\frac{80}{\sin81^\circ}
  \]
  なので
  \[
  AT=\frac{80\sin57^\circ}{\sin81^\circ}\approx67.9\text{ m}.
  \]
  \item
  \[
  PT=AT\tan28^\circ
  =\frac{80\sin57^\circ}{\sin81^\circ}\tan28^\circ
  \approx36.1\text{ m}.
  \]
  \item 仰角が $27^\circ$ のとき
  \[
  PT' = AT\tan27^\circ\approx34.6\text{ m}.
  \]
  よって変化量は
  \[
  PT'-PT=AT(\tan27^\circ-\tan28^\circ)\approx-1.5\text{ m}.
  \]
  したがって、塔の高さは約 $1.5\text{ m}$ 減少する。
\end{enumerate}

\section*{第2問（円）}
\begin{enumerate}[label=\arabic*.]
  \item 弦の公式 $AB=2R\sin\dfrac\theta2$（$\theta=\angle AOB$）より
  \[
  6\sqrt3=12\sin\frac\theta2
  \Rightarrow \sin\frac\theta2=\frac{\sqrt3}{2}.
  \]
  小さい方より
  \[
  \theta=120^\circ.
  \]
  \item 弧長
  \[
  \frac{120}{360}\cdot2\pi\cdot6=4\pi.
  \]
  扇形面積
  \[
  \frac{120}{360}\cdot\pi\cdot6^2=12\pi.
  \]
  \item 接線と弦のなす角は、対応する円周角に等しい。したがって
  \[
  \angle PAB=\frac12\angle AOB=60^\circ.
  \]
  \item
  \[
  [\triangle PAB]=\frac12\cdot PA\cdot AB\cdot\sin\angle PAB
  =\frac12\cdot6\cdot6\sqrt3\cdot\sin60^\circ=27.
  \]
  \item 余弦定理より
  \[
  PB^2=6^2+(6\sqrt3)^2-2\cdot6\cdot6\sqrt3\cdot\cos60^\circ
  =144-36\sqrt3.
  \]
  よって
  \[
  PB=6\sqrt{4-\sqrt3}.
  \]
\end{enumerate}

\section*{第3問（平面図形）}
\begin{enumerate}[label=\arabic*.]
  \item
  \[
  S=AB\cdot AD\sin60^\circ=8\cdot5\cdot\frac{\sqrt3}{2}=20\sqrt3.
  \]
  \item 三角形 $ABD$ に余弦定理を用いて
  \[
  BD^2=8^2+5^2-2\cdot8\cdot5\cos60^\circ=49.
  \]
  よって
  \[
  BD=7.
  \]
  \item 対角線で4等分されるので
  \[
  [\triangle AOD]=\frac14\cdot20\sqrt3=5\sqrt3.
  \]
  \item
  \[
  [\triangle ABD]=\frac12\cdot20\sqrt3=10\sqrt3.
  \]
  点 $A$ から直線 $BD$ への距離を $AH$ とすると
  \[
  10\sqrt3=\frac12\cdot7\cdot AH
  \Rightarrow AH=\frac{20\sqrt3}{7}.
  \]
  \item 三角形 $ABD$ で余弦定理の形を用いると
  \[
  BH=\frac{AB^2+BD^2-AD^2}{2BD}=\frac{64+49-25}{14}=\frac{44}{7},
  \]
  \[
  HD=\frac{AD^2+BD^2-AB^2}{2BD}=\frac{25+49-64}{14}=\frac{5}{7}.
  \]
  よって
  \[
  BH:HD=44:5.
  \]
\end{enumerate}

\section*{第4問（空間図形）}
座標を
\[
A(0,0,0),\ B(6,0,0),\ C(6,4,0),\ D(0,4,0),\ E(0,0,8),\ G(6,4,8)
\]
とする。
\begin{enumerate}[label=\arabic*.]
  \item
  \[
  M=\left(6,4,4\right),\quad N=\left(0,0,2\right).
  \]
  \item
  \[
  \overrightarrow{BM}=(0,4,4),\quad \overrightarrow{BN}=(-6,0,2).
  \]
  法線ベクトルは
  \[
  \overrightarrow{BM}\times\overrightarrow{BN}=(8,-24,24)\propto(1,-3,3).
  \]
  よって平面方程式は
  \[
  (x-6)-3y+3z=0
  \iff x-3y+3z-6=0.
  \]
  \item 点 $D(0,4,0)$ から平面までの距離は
  \[
  \frac{|0-3\cdot4+0-6|}{\sqrt{1^2+(-3)^2+3^2}}
  =\frac{18}{\sqrt{19}}.
  \]
  \item 線分 $AG$ を
  \[
  (x,y,z)=(6t,4t,8t)\quad(0\le t\le1)
  \]
  とおく。平面式に代入すると
  \[
  6t-12t+24t-6=0
  \Rightarrow 18t=6
  \Rightarrow t=\frac13.
  \]
  よって
  \[
  AX:XG=\frac13:\frac23=1:2.
  \]
\end{enumerate}

\end{document}
